\chapter{Guía de Uso del Sistema}
\label{apendice:guia_uso}

Este apéndice proporciona instrucciones detalladas para instalar, configurar y utilizar la reimplementación del sistema de detección de semáforos.

\section{Requisitos del Sistema}

\subsection{Hardware Recomendado}

\begin{table}[H]
\centering
\begin{tabular}{|l|l|}
\hline
\textbf{Componente} & \textbf{Especificación} \\
\hline
CPU & Intel i5/i7 o AMD equivalente (multi-core) \\
RAM & Mínimo 8 GB, recomendado 16 GB \\
GPU & NVIDIA con soporte CUDA (opcional pero recomendado) \\
Almacenamiento & 5 GB libres \\
\hline
\end{tabular}
\caption{Requisitos de hardware}
\end{table}

\subsection{Software Requerido}

\begin{itemize}
    \item Python 3.7 o superior
    \item CUDA Toolkit (opcional, si se usa GPU)
    \item Sistema operativo: Linux (Ubuntu 18.04+), macOS, o Windows 10+
\end{itemize}

\section{Instalación}

\subsection{Clonar el Repositorio}

\begin{codesnippet}
\begin{verbatim}
git clone [URL_DEL_REPOSITORIO]
cd TrafficLightDetection
\end{verbatim}
\end{codesnippet}

\subsection{Crear Entorno Virtual (Recomendado)}

\begin{codesnippet}
\begin{verbatim}
# Crear entorno virtual
python3 -m venv venv

# Activar entorno
# En Linux/macOS:
source venv/bin/activate
# En Windows:
venv\Scripts\activate
\end{verbatim}
\end{codesnippet}

\subsection{Instalar Dependencias}

\begin{codesnippet}
\begin{verbatim}
pip install -r requirements.txt
\end{verbatim}
\end{codesnippet}

Las principales dependencias incluyen:
\begin{itemize}
    \item PyTorch $\geq$ 1.8.0
    \item torchvision
    \item NumPy
    \item OpenCV (cv2)
    \item Matplotlib
\end{itemize}

\subsection{Verificar Instalación}

\begin{codesnippet}
\begin{verbatim}
# Verificar PyTorch
python -c "import torch; print(torch.__version__)"

# Verificar CUDA (si aplica)
python -c "import torch; print(torch.cuda.is_available())"
\end{verbatim}
\end{codesnippet}

\section{Uso Básico}

\subsection{API Principal}

El punto de entrada principal es la función \texttt{load\_pipeline()}:

\begin{codesnippet}
\begin{verbatim}
from tlr.pipeline import load_pipeline
import cv2

# Cargar pipeline
# Para GPU:
pipeline = load_pipeline('cuda:0')
# Para CPU:
# pipeline = load_pipeline('cpu')

# Cargar imagen
image = cv2.imread('path/to/image.jpg')

# Definir projection boxes
# Formato: [x1, y1, x2, y2] para cada ROI
projection_boxes = [
    [100, 50, 150, 120],  # box 1
    [500, 80, 550, 150],  # box 2
]

# Ejecutar detección
result = pipeline(image, projection_boxes)

# Resultado contiene:
# - result['bboxes']: coordenadas de semáforos detectados
# - result['colors']: estado de cada semáforo
# - result['scores']: confianza de cada detección
\end{verbatim}
\end{codesnippet}

\subsection{Procesamiento de Video}

\begin{codesnippet}
\begin{verbatim}
import cv2
from tlr.pipeline import load_pipeline

# Cargar pipeline
pipeline = load_pipeline('cuda:0')

# Abrir video
cap = cv2.VideoCapture('video.mp4')
frame_count = 0

while cap.isOpened():
    ret, frame = cap.read()
    if not ret:
        break

    # Procesar frame
    # (projection_boxes debería venir de HD-map o archivo)
    result = pipeline(frame, projection_boxes)

    # Visualizar resultados
    for bbox, color, score in zip(result['bboxes'],
                                   result['colors'],
                                   result['scores']):
        x1, y1, x2, y2 = bbox
        cv2.rectangle(frame, (x1, y1), (x2, y2), (0, 255, 0), 2)
        cv2.putText(frame, f'{color} ({score:.2f})',
                   (x1, y1-10), cv2.FONT_HERSHEY_SIMPLEX,
                   0.5, (0, 255, 0), 2)

    # Mostrar frame
    cv2.imshow('Traffic Light Detection', frame)
    if cv2.waitKey(1) & 0xFF == ord('q'):
        break

    frame_count += 1

cap.release()
cv2.destroyAllWindows()
\end{verbatim}
\end{codesnippet}

\section{Scripts de Ejemplo}

\subsection{Detección en Imagen Individual}

\begin{codesnippet}
\begin{verbatim}
python example.py
\end{verbatim}
\end{codesnippet}

Procesa una imagen de ejemplo y muestra resultados.

\subsection{Detección con Tracking}

\begin{codesnippet}
\begin{verbatim}
python example_with_tracking.py
\end{verbatim}
\end{codesnippet}

Demuestra el uso del módulo de tracking temporal.

\subsection{Procesamiento en Batch}

\begin{codesnippet}
\begin{verbatim}
python run_pipeline_batch.py
\end{verbatim}
\end{codesnippet}

Procesa múltiples imágenes de forma eficiente.

\subsection{Modo Debug}

\begin{codesnippet}
\begin{verbatim}
python run_pipeline_debug.py
\end{verbatim}
\end{codesnippet}

Ejecuta el pipeline con salida detallada de debugging y visualizaciones intermedias.

\section{Herramientas Auxiliares}

\subsection{Extracción de Frames de Video}

\textbf{Manual:}
\begin{codesnippet}
\begin{verbatim}
python extract_frames_manual.py --video input.mp4 --output frames/
\end{verbatim}
\end{codesnippet}

\textbf{Automática:}
\begin{codesnippet}
\begin{verbatim}
python auto_extraction_of_frames.py --video input.mp4
\end{verbatim}
\end{codesnippet}

\subsection{Selección de Projection Boxes}

\textbf{Interactiva:}
\begin{codesnippet}
\begin{verbatim}
python select_projection_and_append.py --image frame.jpg
\end{verbatim}
\end{codesnippet}

Permite seleccionar ROIs manualmente con el mouse.

\textbf{Programática:}
\begin{codesnippet}
\begin{verbatim}
python projection_boxes_generator.py --config config.json
\end{verbatim}
\end{codesnippet}

\section{Ejecución de Tests}

\subsection{Suite Completa de Tests}

\begin{codesnippet}
\begin{verbatim}
# Ejecutar todos los tests
pytest tests/

# Con salida verbose
pytest -v tests/

# Con medición de cobertura
pytest --cov=src/tlr tests/
\end{verbatim}
\end{codesnippet}

\subsection{Tests por Categoría}

\begin{codesnippet}
\begin{verbatim}
# Solo tests de validación
pytest tests/validation/

# Solo tests exploratorios
pytest tests/exploratory/

# Solo tests de propiedades
pytest tests/properties/

# Solo tests metamórficos
pytest tests/metamorphic/
\end{verbatim}
\end{codesnippet}

\subsection{Tests Específicos}

\begin{codesnippet}
\begin{verbatim}
# Test específico de módulo
pytest tests/validation/test_detection.py

# Test específico de función
pytest tests/validation/test_detection.py::test_nms_basic
\end{verbatim}
\end{codesnippet}

\section{Configuración}

\subsection{Archivos de Configuración}

Los modelos usan archivos de configuración en \texttt{src/tlr/confs/}:

\begin{verbatim}
src/tlr/confs/
├── detector_config.json     # Configuración del detector
├── hori_config.json         # Clasificador horizontal
├── vert_config.json         # Clasificador vertical
└── quad_config.json         # Clasificador quad
\end{verbatim}

\subsection{Modificar Parámetros}

Ejemplo de modificación de parámetros:

\begin{codesnippet}
\begin{verbatim}
{
  "crop_scale": 2.5,
  "min_crop_size": 270,
  "nms_threshold": 0.6,
  "confidence_threshold": 0.3,
  "means_detector": [102.98, 115.95, 122.77],
  "means_recognizer": [66.56, 66.58, 69.06]
}
\end{verbatim}
\end{codesnippet}

\section{Solución de Problemas}

\subsection{CUDA Out of Memory}

Si encuentra errores de memoria en GPU:

\begin{itemize}
    \item Reducir el tamaño de batch (si procesa múltiples imágenes)
    \item Procesar imágenes de menor resolución
    \item Usar CPU en lugar de GPU: \texttt{load\_pipeline('cpu')}
\end{itemize}

\subsection{Rendimiento Bajo en CPU}

El sistema está optimizado para GPU. En CPU:

\begin{itemize}
    \item Espere tiempos de procesamiento más lentos (5-10x)
    \item Considere procesar en batch para amortizar overhead
    \item Reduzca resolución de entrada si es posible
\end{itemize}

\subsection{Errores de OpenCV}

Si hay conflictos con OpenCV:

\begin{codesnippet}
\begin{verbatim}
pip uninstall opencv-python opencv-contrib-python
pip install opencv-python==4.5.5.64
\end{verbatim}
\end{codesnippet}

\subsection{Modelos No Encontrados}

Si falta algún archivo de modelo:

\begin{codesnippet}
\begin{verbatim}
# Verificar que existen los modelos
ls src/tlr/weights/

# Deberían estar:
# - tl.torch
# - hori.torch
# - vert.torch
# - quad.torch
\end{verbatim}
\end{codesnippet}

\section{Estructura de Directorios del Proyecto}

\begin{verbatim}
TrafficLightDetection/
├── src/
│   └── tlr/
│       ├── pipeline.py           # Pipeline principal
│       ├── detector.py           # Detector de semáforos
│       ├── recognizer.py         # Clasificadores
│       ├── tracking.py           # Tracking temporal
│       ├── selector.py           # Asignación húngara
│       ├── hungarian_optimizer.py
│       ├── utils.py              # Utilidades
│       ├── weights/              # Modelos pre-entrenados
│       └── confs/                # Configuraciones
├── tests/
│   ├── validation/               # Tests oracle-based
│   ├── properties/               # Tests basados en propiedades
│   ├── exploratory/              # Tests exploratorios
│   └── metamorphic/              # Tests metamórficos
├── inputTesis/                   # Videos de entrada
├── frames_labeled/               # Frames con ground truth
├── frames_auto_labeled/          # Frames auto-etiquetados
├── examples/                     # Outputs de ejemplo
├── docs/                         # Documentación
│   ├── tesis/                    # Documentos de la tesis
│   └── ...
├── example.py                    # Script de ejemplo básico
├── example_with_tracking.py      # Ejemplo con tracking
├── run_pipeline_batch.py         # Procesamiento en batch
├── run_pipeline_debug.py         # Modo debug
├── requirements.txt              # Dependencias
└── README.md                     # Readme del proyecto
\end{verbatim}

\section{Contacto y Soporte}

Para preguntas, issues o contribuciones:

\begin{itemize}
    \item \textbf{Repositorio:} [URL DEL REPOSITORIO]
    \item \textbf{Issues:} [URL DE ISSUES]
    \item \textbf{Email:} cirob100@gmail.com
\end{itemize}